\clearpage
\selectlanguage{english}
\begin{thebibliography}{99}
\bibitem{1}
Albert Einstein. Die Grundlage der allgemeinen Relativitätstheorie. Annalen der Physik 354, 769-822 (1916).
\par DOI: \nolinkurl{10.1002/andp.19163540702}.
\par\url{https://doi.org/10.1002/andp.19163540702}
\bibitem{2}
Werner Heisenberg. Über den anschaulichen Inhalt der quantentheoretischen Kinematik und Mechanik. Zeitschrift für Physik 43, 172-198 (1927).
\par DOI: \nolinkurl{10.1007/BF01397280}.
\par\url{https://doi.org/10.1007/BF01397280}
\bibitem{3}
H. P. Robertson. The Uncertainty Principle. Physical Review 34, 163-164 (1929).
\par DOI: \nolinkurl{10.1103/PhysRev.34.163}.
\par\url{https://doi.org/10.1103/PhysRev.34.163}
\bibitem{4}
Erwin Schrödinger. Die gegenwärtige Situation in der Quantenmechanik. Die Naturwissenschaften 23, 807-812 (1935), Part I, especially Sec. 5.
\par DOI: \nolinkurl{10.1007/BF01491891}.
\par\url{https://doi.org/10.1007/BF01491891}
\bibitem{5}
Wojciech H. Zurek. Decoherence, einselection, and the quantum origins of the classical. Reviews of Modern Physics 75, 715-775 (2003), arXiv version v3.
\par DOI: \nolinkurl{10.1103/RevModPhys.75.715}.
\par\url{https://arxiv.org/abs/quant-ph/0105127v3}
\par arXiv:quant-ph/0105127v3.
\bibitem{6}
Charles H. Bennett. Logical Reversibility of Computation. IBM Journal of Research and Development 17, 525-532 (1973).
\par DOI: \nolinkurl{10.1147/rd.176.0525}.
\par\url{https://doi.org/10.1147/rd.176.0525}
\bibitem{7}
Ming Li and Paul Vitányi. Reversibility and Adiabatic Computation: Trading Time and Space for Energy. arXiv:quant-ph/9703022 (1997).
\par DOI: \nolinkurl{10.48550/arXiv.quant-ph/9703022}.
\par\url{https://arxiv.org/abs/quant-ph/9703022}
\par arXiv:quant-ph/9703022.
\bibitem{8}
Claude E. Shannon. A Mathematical Theory of Communication. Bell System Technical Journal 27, 379-423 (1948), Part I.
\par DOI: \nolinkurl{10.1002/j.1538-7305.1948.tb01338.x}.
\par\url{https://doi.org/10.1002/j.1538-7305.1948.tb01338.x}
\bibitem{9}
Rolf Landauer. Irreversibility and Heat Generation in the Computing Process. IBM Journal of Research and Development 5, 183-191 (1961).
\par DOI: \nolinkurl{10.1147/rd.53.0183}.
\par\url{https://doi.org/10.1147/rd.53.0183}
\bibitem{10}
Sabine Hossenfelder and Tim Palmer. Rethinking Superdeterminism. Frontiers in Physics 8, Article 139 (2020), especially Secs. 3-4.
\par DOI: \nolinkurl{10.3389/fphy.2020.00139}.
\par\url{https://www.frontiersin.org/journals/physics/articles/10.3389/fphy.2020.00139/full}
\bibitem{11}
B. Hensen et al.. Experimental loophole-free violation of a Bell inequality using entangled electron spins separated by 1.3 km. Nature 526, 682-686 (2015), arXiv:1508.05949v1.
\par DOI: \nolinkurl{10.1038/nature15759}.
\par\url{https://arxiv.org/abs/1508.05949v1}
\par arXiv:1508.05949v1.
\bibitem{12}
Shinsei Ryu and Tadashi Takayanagi. Holographic Derivation of Entanglement Entropy from AdS/CFT. Physical Review Letters 96, 181602 (2006), arXiv version v2.
\par DOI: \nolinkurl{10.1103/PhysRevLett.96.181602}.
\par\url{https://arxiv.org/abs/hep-th/0603001v2}
\par arXiv:hep-th/0603001v2.
\bibitem{13}
Mark Van Raamsdonk. Building up spacetime with quantum entanglement. General Relativity and Gravitation 42, 2323-2329 (2010), original arXiv version v1.
\par DOI: \nolinkurl{10.1007/s10714-010-1034-0}.
\par\url{https://arxiv.org/abs/1005.3035v1}
\par arXiv:1005.3035v1.
\bibitem{14}
Luca Bombelli, Joohan Lee, David Meyer and Rafael D. Sorkin. Space-time as a causal set. Physical Review Letters 59, 521-524 (1987).
\par DOI: \nolinkurl{10.1103/PhysRevLett.59.521}.
\par\url{https://doi.org/10.1103/PhysRevLett.59.521}
\bibitem{15}
K. B. Wharton and N. Argaman. Bell's Theorem and Locally-Mediated Reformulations of Quantum Mechanics. Reviews of Modern Physics 92, 021002 (2020), arXiv version v3.
\par DOI: \nolinkurl{10.1103/RevModPhys.92.021002}.
\par\url{https://arxiv.org/abs/1906.04313v3}
\par arXiv:1906.04313v3.
\bibitem{16}
Maximilian Schlosshauer. Decoherence, the measurement problem, and interpretations of quantum mechanics. Reviews of Modern Physics 76, 1267-1305 (2004), final arXiv version v4 of 2005, especially Secs. II-III.
\par DOI: \nolinkurl{10.1103/RevModPhys.76.1267}.
\par\url{https://arxiv.org/abs/quant-ph/0312059v4}
\par arXiv:quant-ph/0312059v4.
\bibitem{17}
Joel L. Lebowitz. Boltzmann's Entropy and Time's Arrow. Physics Today 46(9), 32-38 (1993).
\par DOI: \nolinkurl{10.1063/1.881363}.
\par\url{https://physicstoday.aip.org/features/boltzmanns-entropy-and-times-arrow}
\bibitem{18}
Charles H. Bennett. Notes on Landauer's principle, reversible computation, and Maxwell's Demon. Studies in History and Philosophy of Modern Physics 34, 501-510 (2003), arXiv:physics/0210005v2.
\par DOI: \nolinkurl{10.1016/S1355-2198(03)00039-X}.
\par\url{https://arxiv.org/abs/physics/0210005v2}
\par arXiv:physics/0210005v2.
\bibitem[Q1]{Q1}
Ingolf Lohmann. Von Softwarearchitektur zur Weltformel - DAS UNIVERSUM ALS ROUND TRIP. Version 1.0.0-candidate.1 (2026), especially Secs. 1, 3, 4 and 6.
\par DOI: \nolinkurl{10.5281/zenodo.21888130}.
\par\url{https://zenodo.org/records/21888130}
\end{thebibliography}
